% ---
% id: transmitter/decoy/rem:cutoffs
% kind: rem
% title: "Entropy, source and receiver cutoffs"
% summary: >-
%   The entropy cutoff bounds the quantum sectors and the source cutoff the
%   flagged ones, and raising the second never requires raising the first. The
%   receiver's cutoff is a different object, which no condition of this
%   submodule reads.
% prominence: standard
% status: accepted
% uses: [def:tagged-source]
% mentions: [cor:emitted-decoy, req:entropy-cutoff]
% assumes: []
% refs: []
% provenance: newly-derived
% lean:
% notes: "The receiver cutoff is named only to be excluded: no claim of this submodule reads it, and the two sides meet only through the compiler."
% ---

\begin{remark}\label{rem:cutoffs}
A reader of a protocol with a decoy source meets three cutoffs that are easily
conflated, and only two of them are data of this submodule.

\emph{The entropy cutoff $N_{\mathrm{ent}}$} is the largest photon number whose
sector of $\Avirt$ stays quantum: at $\mathsf{n}\le N_{\mathrm{ent}}$ the sector
is $S_{\mathsf{n}}$ and carries the signal states $\ket{\mathsf{n},a}$
themselves, one sector for all polarisations of that photon number. Above it the
polarisation is tagged: the sector is one dimensional and there is one per
polarisation, so the declaration separates what a quantum sector keeps
together. Whether keeping a photon number quantum or tagging it per
polarisation gives the higher key rate is a question for the consuming program,
settled by computation and not here.

\emph{The source cutoff $N_{\mathrm{src}}$} is the largest photon number that
keeps a sector of its own at all. Raising it by one adds at most
$|\mathrm{Pol}|$ one-dimensional sectors, one per polarisation some letter
carries at a positive weight, the index set of
Definition~\ref{def:tagged-source}\,(T3) being read off the weights; and it
moves the weight $\tilde p(\mathsf{n}\mid\mu)$ out of the tail flag into them,
so that the weight of that photon number sits on a flag shared by every
intensity instead of on a flag carried by one intensity alone. Raising
$N_{\mathrm{src}}$ never requires raising $N_{\mathrm{ent}}$: (T1) demands only
$N_{\mathrm{ent}}\le N_{\mathrm{src}}$, and the sectors between the two cutoffs
are one dimensional whatever $N_{\mathrm{src}}$ is.

\emph{The receiver cutoff}, the photon number at which a receiver-side
reduction stops retaining blocks and starts flagging, is a datum of the
receiver and of the reduction that consumes it, and not of this submodule. No
condition of this submodule mentions it, no claim here reads it, and nothing
ties it to the entropy or the source cutoff: the transmitter's data and the
receiver's meet only where the protocol compiler composes them. A declaration
with $N_{\mathrm{ent}}=1$ and a receiver cutoff of two is as well formed as any
other.

Every claim of this submodule is stated for arbitrary
$N_{\mathrm{ent}}\le N_{\mathrm{src}}$. The bound $N_{\mathrm{ent}}\le2$ is a
requirement of the emission, the step that outputs the declaration
(Requirement~\ref{req:entropy-cutoff}), and concerns which numbers a
declaration can carry exactly, not any of the mathematics here: for the four
polarisations $\Dset$, at one and at two photons,
Corollary~\ref{cor:emitted-decoy}\,(iv)(b) exhibits a declared basis in which
the signal dyads have the entries $0$ and $\pm\tfrac12$ (and at one photon
also $1$); at three photons and above no such basis is given here.
\end{remark}
